\documentclass[aspectratio=169]{beamer}

\usetheme{Madrid}
\setbeamertemplate{navigation symbols}{}
\setbeamertemplate{footline}[frame number]

\title{Midterm Progress Presentation}
\subtitle{OpenClaw-Driven Sagittarius Arm System}
\author{Your Team Name}
\date{April 14, 2026}

\begin{document}

\begin{frame}
  \titlepage
\end{frame}

\begin{frame}{Motivation}
  \begin{columns}[T]
    \begin{column}{0.48\textwidth}
      \begin{block}{Why this route?}
        \begin{itemize}
          \item Recent 2025--2026 work suggests that purely open-loop LLM control is fragile for real robots, especially in long-horizon and dynamic tasks.
          \item Agentic robotics papers emphasize structured skills, closed-loop feedback, and reflection as practical ways to improve reliability.
          \item For a small team, OpenClaw + Sagittarius gives us a realistic path from natural language to grounded manipulation without training a large robot policy from scratch.
        \end{itemize}
      \end{block}
    \end{column}
    \begin{column}{0.48\textwidth}
      \begin{block}{Why is it useful?}
        \begin{itemize}
          \item Lower data and training cost than an end-to-end VLA route
          \item Easier grounding, debugging, safety control, and failure recovery
          \item A reusable skill interface that can still support future VLA-style integration
        \end{itemize}
      \end{block}
    \end{column}
  \end{columns}

  \vspace{0.5em}

  \begin{block}{Research Motivation}
    Inspired by recent work on agentic robotics and embodied reasoning, we treat the robot as a system of reusable skills controlled by an AI agent, rather than relying only on manual scripts or a fully end-to-end policy.
  \end{block}
\end{frame}

\begin{frame}{Goal}
  \begin{block}{Project Goal}
    Build a first-stage embodied agent system in which OpenClaw converts natural-language tasks into reusable robotic skills on the Sagittarius arm.
  \end{block}

  \vspace{0.5em}

  \begin{itemize}
    \item Provide a unified task interface for search, detection, pick, place, and sorting.
    \item Keep the execution grounded through perception and feedback-aware robot control.
    \item Show that modular agent-driven robotics is a practical and explainable route for real-world task execution.
    \item Establish an extensible system foundation for future long-horizon behaviors and possible VLA integration.
  \end{itemize}
\end{frame}

\begin{frame}{Methodology}
  \begin{block}{Closed-Loop Task Execution Pipeline}
    \centering
    Natural-language task $\rightarrow$ Agent reasoning and skill selection $\rightarrow$ Unified robotic skill call $\rightarrow$ Perception-guided motion execution $\rightarrow$ Feedback to the agent
  \end{block}

  \vspace{0.5em}

  \begin{itemize}
    \item OpenClaw interprets user intent and selects a high-level robot skill.
    \item The skill interface constrains the action space to reusable operations such as search, detect, pick, place, and sort.
    \item The execution layer grounds each skill through visual perception, target localization, motion planning, and arm control.
    \item Returned execution feedback supports safer, more interpretable, and more extensible agent behavior.
  \end{itemize}
\end{frame}

\begin{frame}{Why This Route Before End-to-End VLA?}
  \begin{itemize}
    \item Compared with an end-to-end VLA route, our current framework requires much less task-specific data and training cost.
    \item The system is easier to interpret and debug because perception, decision, and execution are separated into clear layers.
    \item It is more practical for a small team and limited hardware setup, since we can reuse mature robotics modules instead of training a large policy from scratch.
    \item Safety and evaluation are also easier, because the agent calls constrained skills rather than directly generating all low-level actions.
  \end{itemize}

  \vspace{0.5em}

  \begin{block}{Our Position}
    This is not an alternative to VLA in principle. It is a practical intermediate route that builds the infrastructure, interfaces, and task abstractions that can later support VLA-style methods as well.
  \end{block}
\end{frame}

\begin{frame}{Progress Achieved}
  \begin{itemize}
    \item We have encapsulated the core robotic arm capabilities into reusable commands instead of relying on scattered demo scripts.
    \item We have established a complete control path from OpenClaw to the Sagittarius arm.
    \item We have implemented a set of agent-friendly commands, including:
    \begin{itemize}
      \item status, search, detect\_color
      \item pick\_once, pick\_any, pick\_and\_place
      \item classify\_once\_fixed, sort\_all\_fixed, classify\_once\_map
    \end{itemize}
    \item We have also created an OpenClaw skill so the agent can call these capabilities in a consistent way.
  \end{itemize}
\end{frame}

\begin{frame}{What We Can Demonstrate Now}
  \begin{itemize}
    \item System readiness check before task execution
    \item Observation pose and color-based target detection
    \item Single-object grasping for a specified color
    \item Closed-loop pick-and-place for a visible target
    \item Repeated sorting of visible blocks with fixed placement rules
  \end{itemize}

  \vspace{0.5em}

  \begin{block}{Current Value}
    The system has moved from low-level arm control toward task-level robotic capabilities that are suitable for AI agent invocation.
  \end{block}
\end{frame}

\begin{frame}{Why This Matters}
  \begin{itemize}
    \item The project is no longer only about moving a robotic arm; it is about giving an AI agent a structured way to use robotic capabilities.
    \item The abstraction level is higher: the agent decides what to do, while the robotic stack handles how to execute it.
    \item This combination is useful for research on embodied agents, human-robot interaction, and task-level automation.
    \item The current architecture is also a practical foundation for future work in remote invocation, multi-step task execution, and more autonomous decision-making.
  \end{itemize}
\end{frame}

\begin{frame}{Next Steps}
  \begin{itemize}
    \item Improve the robustness of the control pipeline and failure recovery logic.
    \item Extend the interface for remote invocation, such as HTTP or curl-based control.
    \item Strengthen multi-step agent behavior beyond single-command execution.
    \item Prepare more complete demonstrations that better highlight AI-agent advantages.
  \end{itemize}

  \vspace{0.5em}

  \begin{block}{Takeaway}
    At the midterm stage, our main achievement is not only making the arm move, but organizing the robotic capabilities into an interface that an AI agent can meaningfully use.
  \end{block}
\end{frame}

\end{document}
